\documentclass{sajzk}

\usepackage{array}

\usepackage{minted}
% \usepackage{listings}
% \lstset{basicstyle=\ttfamily}
% \usepackage{xcolor}
% \usepackage{mdframed}
% 
% \definecolor{light-gray}{gray}{0.95} % the shade of grey that stack exchange uses

\usetikzlibrary {arrows.meta, angles, quotes}
\usetikzlibrary {calc}

\begin{document}
\section{Axiomatische sichtweise}
\label{il6v}
\lhead{:mathe:ga:}

Es gibt mehrere Axiomensysteme. Hier eins von Hestenes und eins von Chisolm.
\begin{itemize}
    \item \nameref{ciuz} ciuz.tex
    \item \nameref{ompz} ompz.tex
\end{itemize}

\subsection{Referencen}
\begin{itemize}
    \item \nameref{f35d} f35d.tex
\end{itemize}

\subsection{Literatur}
\begin{itemize}
    \item David Hestenes - New Foundations for Classical Mechanics (2002) Seite 35
    \item David Hestenes, Garret Sobczyk - Clifford Algebra to Geometric Calculus A
    \item Unified Language for Mathematics and Physics (1984) Seite 3
    \item Eric Chisolm - Geomeric Algebra (2012) Seite 10-12
\end{itemize}
\end{document}
